% C4_brief.tex -- ONE source for the English and Romanian case brief of Case 4 (Lectures 10 and 13).
\input{numbers_\lang}
\renewcommand{\emqacase}{\T{Research case 4 (Lectures 10 and 13)}{Studiul de caz 4 (cursurile 10 și 13)}}
\begin{document}
\emqatitle{\T{Case 4: A battery on the day-ahead market --- quantile networks and reinforcement learning}%
                {Studiul de caz 4: O baterie pe piața pentru ziua următoare --- rețele neuronale cuantilice și reinforcement learning}}%
          {\T{Based on}{Studiu bazat pe lucrarea} Andrei, Cepoi, Dumitrescu, Lessmann, Chiriac (Luchian), Pele \& Petukhina (2026), \emph{Joint probabilistic forecasting for battery arbitrage in day-ahead power markets},
           SSRN 6872381, \href{https://doi.org/10.2139/ssrn.6872381}{doi:10.2139/ssrn.6872381}}

\begin{scenario}
\T{An energy trader in Bucharest runs a 1 MW / 1 MWh battery energy storage system (BESS). Every day before noon it must submit tomorrow's bids to the day-ahead market without knowing tomorrow's prices.
Two teams pitch a tool. The data-science team proposes \textbf{a neural network that forecasts the whole distribution of tomorrow's 24 prices}, followed by a simple rule that picks the hours.
The quant team proposes \textbf{a reinforcement-learning agent that learns by itself when to charge and when to discharge}.
The board asks two questions: \textbf{which tool earns more on unseen data, and does any of them earn enough to repay the battery?}}%
{Un trader de energie din București operează un sistem de stocare a energiei în baterii (BESS, \emph{battery energy storage system}) de 1 MW / 1 MWh. În fiecare zi, înainte de prînz, trebuie să depună ofertele pentru ziua următoare, fără să cunoască prețurile de a doua zi.
Două echipe propun cîte un instrument. Echipa de știința datelor propune \textbf{o rețea neuronală care prognozează întreaga distribuție a celor 24 de prețuri din ziua următoare}, urmată de o regulă simplă care alege orele.
Echipa de analiză cantitativă propune \textbf{un agent de reinforcement learning care învață singur cînd să încarce și cînd să descarce bateria}.
Consiliul de administrație pune două întrebări: \textbf{care instrument cîștigă mai mult pe date noi și cîștigă vreunul suficient pentru a recupera investiția în baterie?}}
\end{scenario}

\section{\T{The reference paper}{Lucrarea de referință}}
\textbf{\T{Question.}{Întrebarea de cercetare.}} \T{How much arbitrage value can a battery capture \emph{ex ante}, when it is operated on probabilistic price forecasts, and does that value justify the investment, in four interconnected Central and Eastern European (CEE) day-ahead markets (Sections~1 and 3)?}%
{Ce parte din valoarea de arbitraj poate capta o baterie \emph{ex ante}, atunci cînd este operată pe baza prognozelor probabilistice ale prețului, și justifică această valoare investiția, pe patru piețe interconectate pentru ziua următoare din Europa Centrală și de Est (secțiunile~1 și 3)?}

\textbf{\T{Data.}{Datele.}} \T{Hourly day-ahead prices of Romania, Bulgaria, Hungary and the Czech Republic, January 2018 -- September 2024 (59,160 hours per country), with 37 predictors: load, generation and their forecasts, lagged prices, Brent, TTF (Title Transfer Facility) gas and carbon prices, weather and calendar (Section~5.1).
Battery: 1 MW / 1 MWh, efficiency 90\,\%, CAPEX (capital expenditure) 500,000 EUR, 15 years, 7\,\% discount rate, at most one cycle per day (Section~4).}%
{Prețurile orare pentru ziua următoare din România, Bulgaria, Ungaria și Cehia, ianuarie 2018 -- septembrie 2024 (59.160 de ore pentru fiecare țară), cu 37 de variabile explicative: consumul, producția și prognozele acestora, prețurile anterioare, prețul petrolului Brent, prețul gazelor naturale TTF (\emph{Title Transfer Facility}) și prețul certificatelor de emisii de CO$_2$, variabile meteorologice și calendaristice (secțiunea~5.1).
Bateria: 1 MW / 1 MWh, randament de 90\,\%, CAPEX (cheltuieli de capital) de 500.000 EUR, durată de viață de 15 ani, rată de actualizare de 7\,\%, cel mult un ciclu pe zi (secțiunea~4).}

\textbf{\T{Method.}{Metoda.}} \T{Distributional Random Forests (DRF) forecast the \emph{joint} distribution of the 96 next-day prices (24 hours $\times$ 4 countries), refitted on the last 365 days every 30 days.
A quantile-based (QB) rule chooses the buy hour and the sell hour from the forecast medians and places limit bids at forecast quantiles; the benchmarks are a fixed pair of hours (NH) and daily perfect foresight (PF-Daily); the investment is judged with net present value (NPV), payback and break-even CAPEX.}%
{Modelele DRF (\emph{Distributional Random Forests}, păduri aleatoare care estimează distribuții întregi) prognozează distribuția \emph{comună} a celor 96 de prețuri din ziua următoare (24 de ore $\times$ 4 țări) și sînt reestimate la fiecare 30 de zile pe ultimele 365 de zile.
O regulă bazată pe cuantile (QB) alege ora de cumpărare și ora de vînzare pe baza medianelor prognozate și plasează oferte cu preț limită la cuantilele prognozei; reperele sînt o pereche fixă de ore (NH) și previziunea perfectă zilnică (PF-Daily); investiția este evaluată prin valoarea actualizată netă (VAN), perioada de recuperare și CAPEX-ul de echilibru.}

\textbf{\T{Main findings of the paper.}{Principalele rezultate ale lucrării.}}
\begin{itemize}
\item \T{At $\alpha=0.8$, QB beats NH in cumulative profit in all four markets, with relative gains of about 13\,\% in Czechia and 7\,\% in Romania (Section~6.2).
Romania, January 2019 -- September 2024: QB 142,671 EUR, NH 133,230 EUR, PF-Daily 188,249 EUR; payback 20.1 years for QB (single-country strategy table for Romania).}%
{Pentru $\alpha=0{,}8$, strategia QB obține un profit cumulat mai mare decît NH pe toate cele patru piețe, cu cîștiguri relative de aproximativ 13\,\% în Cehia și 7\,\% în România (secțiunea~6.2).
România, ianuarie 2019 -- septembrie 2024: QB 142.671 EUR, NH 133.230 EUR, PF-Daily 188.249 EUR; perioada de recuperare a investiției este de 20,1 ani pentru QB (tabelul pentru România).}
\item \T{Under the baseline CAPEX of 500,000 EUR none of the strategies covers the investment cost; the break-even CAPEX of the QB rule lies between 48\,\% and 75\,\% of the baseline (Section~6.2).}%
{La CAPEX-ul de referință de 500.000 EUR, nicio strategie nu acoperă costul investiției; pentru regula QB, CAPEX-ul de echilibru se situează între 48\,\% și 75\,\% din valoarea de referință (secțiunea~6.2).}
\item \T{Joint cross-border dispatch adds about 5\,\% to QB profit in Romania (149.9 vs 142.7 thousand EUR), ``driven by persistent regional price differences'' (Abstract; Table ``Multi-zone joint-dispatch'').
DRF accuracy for Romania: MAE (mean absolute error) 39.63, RMSE (root mean squared error) 49.60, CRPS (continuous ranked probability score) 28.12 EUR/MWh over 70 rolling windows (Table ``DRF forecasting performance'').}%
{Operarea comună, transfrontalieră, crește profitul QB în România cu aproximativ 5\,\% (149,9 față de 142,7 mii EUR), datorită diferențelor regionale persistente de preț (rezumat; tabelul pentru operarea în mai multe zone).
Acuratețea DRF pentru România: eroarea absolută medie (MAE, \emph{mean absolute error}) 39,63, rădăcina erorii pătratice medii (RMSE, \emph{root mean squared error}) 49,60 și scorul probabilistic ordonat continuu (CRPS, \emph{continuous ranked probability score}) 28,12 EUR/MWh, pe 70 de ferestre mobile (tabelul privind performanța prognozelor DRF).}
\end{itemize}
{\footnotesize\color{FooterGray}\T{The same paper is the basis of Case~6 of the course Energy Markets Quantitative Analysis (random forests); here we replace the forest by a neural network and add a reinforcement-learning agent. Related case of this course: explainable early warning of price spikes (Case~2).}%
{Aceeași lucrare stă la baza studiului de caz~6 din cursul Analiza cantitativă a piețelor de energie (cu păduri aleatoare); aici înlocuim pădurea aleatoare cu o rețea neuronală și adăugăm un agent de reinforcement learning. Studiu de caz înrudit din acest curs: avertizarea timpurie explicabilă pentru vîrfurile de preț (studiul de caz~2).}}

\section{\T{The method}{Metoda}}
\begin{intuition}
\T{A battery earns the spread between a cheap hour and a dear hour, minus the energy lost in charging and discharging. Looking back, the best two hours are obvious; looking forward, they are a forecast.
Two ways to decide: \emph{forecast, then optimise} (a network predicts the distribution of tomorrow's prices; a rule picks the hours), or \emph{learn the policy directly} (an agent tries charging and discharging on past days and keeps what paid).}%
{O baterie cîștigă diferența dintre prețul unei ore ieftine și cel al unei ore scumpe, minus energia pierdută la încărcare și descărcare. Retrospectiv, cele mai bune două ore sînt evidente; în momentul deciziei, ele sînt o prognoză.
Există două moduri de a decide: \emph{întîi prognoza, apoi optimizarea} (o rețea prognozează distribuția prețurilor de mîine, iar o regulă alege orele) sau \emph{învățarea directă a politicii} (un agent încearcă încărcarea și descărcarea pe zilele trecute și păstrează ceea ce a adus profit).}
\end{intuition}
\T{\textbf{Profit of a cycle.} With $p_{d,h}$ the price of hour $h$ on day $d$ and efficiency $\eta=0.9$ on each leg, buying in hour $h_1$ and selling in $h_2>h_1$ earns $\pi_d(h_1,h_2)=\eta\,p_{d,h_2}-p_{d,h_1}/\eta$.
Perfect foresight takes $\max_{h_1<h_2}\pi_d$ each day (an ex post upper bound); fixed hours use the pair with the highest total profit in the previous year.

\textbf{Quantile network.} A multilayer perceptron (MLP) maps the information known at gate closure, $\mathbf x_d$ (the 24 prices of $d-1$ and of $d-7$, the mean prices of Bulgaria and Hungary on $d-1$, day of the week, season), through two hidden layers $\mathbf h^{(l)}=\mathrm{ReLU}(\mathbf W^{(l)}\mathbf h^{(l-1)}+\mathbf b^{(l)})$ with 128 units to $24\times9$ outputs $\hat q_{\tau,d,h}$, $\tau\in\{0.1,\dots,0.9\}$.
It is trained with the pinball (quantile) loss $\mathcal L=\frac{1}{24\cdot 9\,n}\sum_{d,h,\tau}\rho_\tau(p_{d,h}-\hat q_{\tau,d,h})$, $\rho_\tau(u)=\max\{\tau u,(\tau-1)u\}$, whose minimiser is the conditional $\tau$-quantile; Adam, weight decay $10^{-4}$, dropout 0.1, early stopping on a later validation window; quantiles are sorted afterwards so that they do not cross.
Quality: $\mathrm{CRPS}\approx\frac{2}{9}\sum_\tau\rho_\tau(p-\hat q_\tau)$ in EUR/MWh (lower is better) and the coverage of the 10\,\%--90\,\% band (target 80\,\%).
\textbf{Rules.} \emph{Median pair}: choose $(h_1,h_2)$ maximising $\eta\hat q_{0.5,h_2}-\hat q_{0.5,h_1}/\eta$, trade only if it is positive, with market orders. \emph{Quantile bids} (the paper's QB idea): bid to buy at $\hat q_{1-q,h_1}$ and to sell at $\hat q_{q,h_2}$, $q=(1-\alpha)/2$, $\alpha=0.8$.

\textbf{Q-learning agent.} A day is an episode of 24 steps. The state is $s_h=(\mathit{soc}_h,\,b_h,\,b^{\max}_{>h},\,b^{\min}_{>h})$: the state of charge (empty/full), the octile of hour $h$ in the forecast ranking of the day, and the highest and lowest octiles still to come. Actions: wait, charge, discharge; reward $r_h$ = the cash flow of the hour.
The action values are updated by $Q(s_h,a_h)\leftarrow Q(s_h,a_h)+\lambda_n\big[r_h+\max_{a'}Q(s_{h+1},a')-Q(s_h,a_h)\big]$ (discount factor 1 within the day, step $\lambda_n=n^{-0.6}$ after $n$ visits), with $\varepsilon$-greedy exploration from 1 down to 0.05; the policy is then $\arg\max_aQ(s,a)$.
\textbf{Investment.} $\bar\pi_{\mathrm{year}}=365\,\bar\pi_d$; $\mathrm{NPV}=-\mathrm{CAPEX}+\bar\pi_{\mathrm{year}}\sum_{t=1}^{15}1.07^{-t}$; payback $=\mathrm{CAPEX}/\bar\pi_{\mathrm{year}}$; NPV $=0$ needs $\bar\pi_{\mathrm{year}}=\Nneed$ EUR.}%
{\textbf{Profitul unui ciclu.} Fie $p_{d,h}$ prețul orei $h$ din ziua $d$ și $\eta=0{,}9$ randamentul fiecărei operații. Cumpărarea în ora $h_1$ și vînzarea în ora $h_2>h_1$ aduc $\pi_d(h_1,h_2)=\eta\,p_{d,h_2}-p_{d,h_1}/\eta$.
Previziunea perfectă alege în fiecare zi $\max_{h_1<h_2}\pi_d$ (o limită superioară, stabilită \emph{ex post}); strategia cu ore fixe folosește perechea cu cel mai mare profit total din anul anterior.

\textbf{Rețeaua cuantilică.} Un perceptron multistrat (MLP, \emph{multilayer perceptron}) transformă informația cunoscută la închiderea porții, $\mathbf x_d$ (cele 24 de prețuri din zilele $d-1$ și $d-7$, prețurile medii din Bulgaria și Ungaria în ziua $d-1$, ziua săptămînii, sezonul), prin două straturi ascunse $\mathbf h^{(l)}=\mathrm{ReLU}(\mathbf W^{(l)}\mathbf h^{(l-1)}+\mathbf b^{(l)})$ cu cîte 128 de neuroni, în $24\times9$ ieșiri $\hat q_{\tau,d,h}$, $\tau\in\{0{,}1,\dots,0{,}9\}$.
Rețeaua este antrenată cu funcția de pierdere \emph{pinball} (cuantilică) $\mathcal L=\frac{1}{24\cdot 9\,n}\sum_{d,h,\tau}\rho_\tau(p_{d,h}-\hat q_{\tau,d,h})$, $\rho_\tau(u)=\max\{\tau u,(\tau-1)u\}$, al cărei minim se atinge în cuantila condiționată de ordin $\tau$; optimizatorul Adam, penalizare a ponderilor $10^{-4}$, \emph{dropout} 0,1, oprire timpurie pe o fereastră de validare ulterioară; cuantilele sînt apoi ordonate, ca să nu se intersecteze.
Calitatea prognozei: $\mathrm{CRPS}\approx\frac{2}{9}\sum_\tau\rho_\tau(p-\hat q_\tau)$, în EUR/MWh (valorile mici sînt preferabile), și acoperirea intervalului 10\,\%--90\,\% (nivelul-țintă este 80\,\%).
\textbf{Regulile.} \emph{Perechea medianei}: se aleg $(h_1,h_2)$ care maximizează $\eta\hat q_{0.5,h_2}-\hat q_{0.5,h_1}/\eta$ și se tranzacționează numai dacă valoarea este pozitivă, cu ordine la prețul pieței. \emph{Ofertele la cuantile} (ideea regulii QB din lucrare): cumpărare la $\hat q_{1-q,h_1}$ și vînzare la $\hat q_{q,h_2}$, cu $q=(1-\alpha)/2$ și $\alpha=0{,}8$.

\textbf{Agentul Q-learning.} O zi este un episod de 24 de pași. Starea este $s_h=(\mathit{soc}_h,\,b_h,\,b^{\max}_{>h},\,b^{\min}_{>h})$: starea de încărcare (goală/plină), octila orei $h$ în ordinea prognozată a zilei și cea mai mare, respectiv cea mai mică octilă din orele rămase. Acțiunile: așteptare, încărcare, descărcare; recompensa $r_h$ este fluxul de numerar al orei.
Valorile acțiunilor se actualizează prin $Q(s_h,a_h)\leftarrow Q(s_h,a_h)+\lambda_n\big[r_h+\max_{a'}Q(s_{h+1},a')-Q(s_h,a_h)\big]$ (factor de actualizare 1 în cadrul zilei, pas $\lambda_n=n^{-0{,}6}$ după $n$ vizite), cu explorare $\varepsilon$-lacomă, de la 1 pînă la 0,05; politica este apoi $\arg\max_aQ(s,a)$.
\textbf{Investiția.} $\bar\pi_{\mathrm{an}}=365\,\bar\pi_d$; $\mathrm{VAN}=-\mathrm{CAPEX}+\bar\pi_{\mathrm{an}}\sum_{t=1}^{15}1{,}07^{-t}$; perioada de recuperare $=\mathrm{CAPEX}/\bar\pi_{\mathrm{an}}$; pentru VAN $=0$ este nevoie de $\bar\pi_{\mathrm{an}}=\Nneed$ EUR.}
\begin{decision}
\T{Judge a trading tool on profit in a later test period, over several seeds, next to fixed hours and perfect foresight; judge the battery on the realistic profit against the \Nneed\ EUR per year needed to repay it.}%
{Evaluați un instrument de tranzacționare după profitul dintr-o perioadă de test ulterioară, pe mai multe seeds, alături de orele fixe și de previziunea perfectă; evaluați bateria după profitul realist, comparat cu cei \Nneed\ EUR pe an necesari pentru recuperarea investiției.}
\end{decision}

\section{\T{Replication on the course data}{Replicarea pe datele cursului}}
\T{Notebook \texttt{C4\_notebook\_EN.ipynb}: hourly day-ahead prices of Romania, Bulgaria and Hungary (energy-charts.info), delivery days in Central European Time (CET) with 24 prices.
Final forecaster trained on 2019 -- August 2024 (\NnTrainTwo\ days), validation September--December 2024; the agent learns on 2023--2024 (\NnAgentDays\ days) from the out-of-sample forecasts of a network trained on 2019--2022; test: January 2025 -- October 2026 (\NnTest\ days), after the paper's sample. Five seeds for every network and every agent.
Classroom implementation: a neural network instead of DRF; only lagged prices and calendar as predictors; Romania only; one cycle per day.
\textbf{Check on the paper's window} (8 January 2019 -- 30 September 2024), where no forecast is needed: our perfect-foresight profit is \NpfW\ EUR (paper 188,249) and the fixed hours \NnhWPair\ CET earn \NnhW\ EUR (paper 133,230).}%
{Notebook-ul \texttt{C4\_notebook\_RO.ipynb} folosește prețurile orare pentru ziua următoare din România, Bulgaria și Ungaria (energy-charts.info), zilele de livrare în ora Europei Centrale (CET, \emph{Central European Time}) cu 24 de prețuri.
Prognoza finală este antrenată pe ianuarie 2019 -- august 2024 (\NnTrainTwo\ de zile), cu validare pe septembrie--decembrie 2024; agentul învață pe 2023--2024 (\NnAgentDays\ de zile) din prognozele, făcute pe date noi, ale unei rețele antrenate pe 2019--2022; perioada de test este ianuarie 2025 -- octombrie 2026 (\NnTest\ de zile), după eșantionul lucrării. Fiecare rețea și fiecare agent sînt estimați cu cinci seeds.
Față de lucrare, implementarea de la curs folosește o rețea neuronală în locul DRF, numai prețurile anterioare și variabilele calendaristice ca variabile explicative, numai România și un singur ciclu pe zi.
\textbf{Verificare pe perioada din lucrare} (8 ianuarie 2019 -- 30 septembrie 2024), unde nu este nevoie de prognoză: profitul nostru în ipoteza previziunii perfecte este de \NpfW\ EUR (lucrarea: 188.249), iar orele fixe \NnhWPair\ CET aduc \NnhW\ EUR (lucrarea: 133.230).}

\begin{figure}[H]
\begin{minipage}[t]{0.49\linewidth}\centering
\includegraphics[width=\linewidth]{f1_daily_shape}
\caption{\T{Mean hourly price shape in Romania, four periods.}{Profilul orar mediu al prețului în România, în patru perioade.}}
\src{energy-charts.info; 2019--2026.}
\end{minipage}\hfill
\begin{minipage}[t]{0.49\linewidth}\centering
\includegraphics[width=\linewidth]{f2_forecast_fan}
\caption{\T{Quantile MLP forecast of one test day (mean of five seeds).}{Prognoza MLP cuantilic pentru o zi de test (media celor cinci seeds).}}
\src{\T{our network; energy-charts.info.}{rețeaua noastră; energy-charts.info.}}
\end{minipage}
\end{figure}

\T{\textbf{Forecast quality on the test period.} The quantile MLP has MAE \Nmae\ EUR/MWh (standard deviation over seeds \NmaeSd), CRPS \Ncrps\ EUR/MWh and a 10\,\%--90\,\% band that covers \Ncover\,\% of hours (target 80\,\%); the naive ``same hour yesterday'' has MAE \NmaeNaive.
The paper's DRF figures (MAE 39.63, CRPS 28.12) refer to 2019--2024, which includes the 2022 crisis, so the levels are not directly comparable.}%
{\textbf{Calitatea prognozei în perioada de test.} MLP-ul cuantilic are MAE \Nmae\ EUR/MWh (abaterea standard între seeds este \NmaeSd), CRPS \Ncrps\ EUR/MWh, iar intervalul 10\,\%--90\,\% cuprinde prețul efectiv în \Ncover\,\% din ore (nivelul-țintă este 80\,\%); prognoza naivă „aceeași oră din ziua precedentă” are MAE \NmaeNaive.
Valorile DRF din lucrare (MAE 39,63, CRPS 28,12) se referă la 2019--2024, perioadă care include criza din 2022, astfel că nivelurile nu sînt direct comparabile.}

\begin{table}[H]\centering\small
\caption{\T{Battery strategies in Romania, 1 MW / 1 MWh, test period January 2025 -- October 2026; networks and agents: mean of five seeds (SD = standard deviation of the profit over seeds).}%
{Strategiile de operare a bateriei în România, 1 MW / 1 MWh, perioada de test ianuarie 2025 -- octombrie 2026; rețele și agenți: media celor cinci seeds (AS = abaterea standard a profitului între seeds).}}
\tabin{strategies}
\src{\T{our computation; CAPEX 500,000 EUR, 15 years, 7\,\% (the paper's assumptions).}{calcule proprii; CAPEX 500.000 EUR, 15 ani, 7\,\% (ipotezele lucrării).}}
\end{table}

\begin{figure}[H]
\begin{minipage}[t]{0.52\linewidth}\centering
\includegraphics[width=\linewidth]{f3_cumulative_profit}
\caption{\T{Cumulative profit on the test period, by strategy.}{Profitul cumulat în perioada de test, pe strategii.}}
\src{\T{our computation; energy-charts.info, 2025--2026.}{calcule proprii; energy-charts.info, 2025--2026.}}
\end{minipage}\hfill
\begin{minipage}[t]{0.46\linewidth}\centering
\includegraphics[width=\linewidth]{f5_learning_curve}
\caption{\T{Learning curve of the Q-learning agents (band: range over five seeds).}{Curba de învățare a agenților Q-learning (banda: intervalul celor cinci seeds).}}
\src{\T{our computation; training days 2023--2024.}{calcule proprii; zilele de antrenare 2023--2024.}}
\end{minipage}
\end{figure}

\T{\textbf{What our data show.} The daily shape changed: the fixed pair chosen on 2024 is \NpairNew\ CET (solar middays are cheap), whereas on the paper's window it was \NnhWPair.
On the test period perfect foresight earns \Npf\ EUR; the fixed hours earn \Nnh\ EUR and the naive forecast \Nnaive.
The quantile MLP with the median pair earns \Nmkt\ EUR (SD over seeds \NmktSd), \NbestShare\,\% of the perfect-foresight profit; with the paper's quantile bids it earns \Nqb\ EUR (both bids accepted on \NqbBoth\,\% of days).
The Q-learning agent learns a sensible policy (Figure~\ref{fig:hours}) and earns \Nrl\ EUR with the MLP signal (SD \NrlSd) and \Nrln\ EUR with yesterday's prices (SD \NrlnSd), about as much as the fixed hours: forecast-then-optimise wins here, because the planner sees the whole forecast curve while the agent sees it through a few ranks, and the agent learnt in 2023--2024, a different daily shape.}%
{\textbf{Ce arată datele.} Profilul zilnic s-a schimbat: perechea fixă aleasă pe 2024 este \NpairNew\ CET (energia de la prînz este ieftină datorită producției solare), în timp ce pe perioada din lucrare era \NnhWPair.
În perioada de test, previziunea perfectă aduce \Npf\ EUR; orele fixe aduc \Nnh\ EUR, iar prognoza naivă \Nnaive.
MLP-ul cuantilic cu perechea medianei cîștigă \Nmkt\ EUR (abaterea standard între seeds este \NmktSd), adică \NbestShare\,\% din profitul previziunii perfecte; cu ofertele la cuantile din lucrare cîștigă \Nqb\ EUR (ambele oferte sînt acceptate în \NqbBoth\,\% din zile).
Agentul Q-learning învață o politică rezonabilă (Figura~\ref{fig:hours}), și cîștigă \Nrl\ EUR cu semnalul MLP (AS \NrlSd) și \Nrln\ EUR cu prețurile de ieri (AS \NrlnSd), aproximativ cît orele fixe: aici cîștigă abordarea „întîi prognoza, apoi optimizarea”, deoarece regula vede întreaga curbă prognozată, iar agentul o vede doar prin cîteva ranguri, și pentru că agentul a învățat pe 2023--2024, cu un alt profil zilnic.}

\begin{figure}[H]\centering
\includegraphics[width=0.78\linewidth]{f4_hours_used}
\caption{\T{When does each strategy charge and discharge? Share of test days by hour (agents and networks: five seeds pooled).}{Cînd încarcă și cînd descarcă fiecare strategie? Proporția zilelor de test pe ore (agenți și rețele: cele cinci seeds cumulate).}}\label{fig:hours}
\src{\T{our computation; January 2025 -- October 2026.}{calcule proprii; ianuarie 2025 -- octombrie 2026.}}
\end{figure}

\T{\textbf{Decision number.} The best realistic tool, the quantile MLP with the median pair, earns \NbestYear\ EUR per year against the \Nneed\ EUR needed: NPV \NbestNPV\ EUR, payback \NbestPay\ years, break-even CAPEX \NbestBe\ EUR.
Even perfect foresight (\NpfYear\ EUR per year, NPV \NpfNPV\ EUR) stays below the threshold: as in the paper, arbitrage alone does not repay a 500,000 EUR battery; a lower CAPEX or other revenues are needed.}%
{\textbf{Cifra relevantă pentru decizie.} Cel mai bun instrument realist, MLP-ul cuantilic cu perechea medianei, cîștigă \NbestYear\ EUR pe an, față de cei \Nneed\ EUR necesari: VAN \NbestNPV\ EUR, perioadă de recuperare de \NbestPay\ ani, CAPEX de echilibru \NbestBe\ EUR.
Chiar și previziunea perfectă (\NpfYear\ EUR pe an, VAN \NpfNPV\ EUR) rămîne sub prag: ca în lucrare, arbitrajul singur nu recuperează o baterie de 500.000 EUR; este nevoie de un CAPEX mai mic sau de alte surse de venit.}

\section{\T{Discussion questions for managers}{Întrebări pentru discuție}}
\begin{enumerate}
\item \T{The quant team says the agent will improve with more data. What evidence (test period, seeds, benchmarks) would you require before giving it the battery?}%
{Echipa de analiză cantitativă susține că agentul se va îmbunătăți cu mai multe date. Ce dovezi (perioadă de test, seeds, repere) ați cere înainte de a-i încredința bateria?}
\item \T{The fixed pair moved from night$\rightarrow$evening to midday$\rightarrow$evening. How often would you retrain the network and the agent, and who decides it?}%
{Perechea fixă s-a mutat de la noapte$\rightarrow$seară la prînz$\rightarrow$seară. Cît de des ați reantrena rețeaua și agentul și cine decide acest lucru?}
\item \T{Why can quantile bids earn less than market orders when the forecast is good, although the paper finds them better? What does the acceptance rate tell you?}%
{De ce pot aduce ofertele la cuantile mai puțin decît ordinele la prețul pieței cînd prognoza este bună, deși lucrarea le găsește mai avantajoase? Ce vă spune rata de acceptare?}
\item \T{Which other revenues (balancing, frequency reserves, capacity payments, a second daily cycle) would you add to close the gap to \Nneed\ EUR per year, and how certain are they?}%
{Ce alte venituri (din piața de echilibrare, din rezervele de reglaj al frecvenței, din plățile pentru capacitate, dintr-un al doilea ciclu zilnic) ați adăuga pentru a atinge \Nneed\ EUR pe an și cît de sigure sînt acestea?}
\item \T{When is reinforcement learning worth its complexity for storage: one battery and one market, or many assets, several markets and degradation costs?}%
{Cînd își merită reinforcement learning complexitatea în cazul stocării: pentru o baterie și o piață sau pentru mai multe active, mai multe piețe și costuri de degradare?}
\end{enumerate}

\section{\T{Limitations of the classroom replication}{Limitele replicării}}
\T{(i) One cycle per day and 1 MWh; real projects use 2--4 hour batteries and several revenue streams. (ii) Romania only, no cross-border dispatch, fees, intraday or balancing markets, no price impact.
(iii) The network uses only lagged prices and calendar variables (the paper's DRF uses 37 predictors) and forecasts the 24 marginal distributions, not their joint distribution.
(iv) The agent's state compresses the forecast into ranks, and it is trained once on 2023--2024; perfect foresight is an ex post bound. (v) No degradation cost and no grid tariffs.}%
{(i) Un singur ciclu pe zi și o capacitate de 1 MWh; proiectele reale folosesc baterii cu durata de descărcare de 2--4 ore și mai multe surse de venit. (ii) Numai România, fără operare transfrontalieră, tarife, piață intrazilnică sau de echilibrare și fără efectul ofertelor asupra prețului.
(iii) Rețeaua folosește numai prețurile anterioare și variabilele calendaristice (modelul DRF din lucrare folosește 37 de variabile explicative) și prognozează cele 24 de distribuții marginale, nu distribuția lor comună.
(iv) Starea agentului rezumă prognoza prin ranguri, iar agentul este antrenat o singură dată, pe 2023--2024; previziunea perfectă este o limită stabilită \emph{ex post}. (v) Nu includem costul degradării bateriei și nici tarifele de rețea.}

\newpage
\section*{\T{Summary of case 4}{Studiul de caz 4 pe scurt}}
\glance{\T{Forecast-then-optimise or a learning agent: which earns more with a 1 MW / 1 MWh battery, and does it repay the battery?}{Prognoză și optimizare sau un agent care învață: care cîștigă mai mult cu o baterie de 1 MW / 1 MWh și se recuperează investiția?}}%
{\T{Hourly day-ahead prices of Romania, Bulgaria and Hungary, 2019--October 2026; test 2025--2026.}{Prețuri orare pentru ziua următoare din România, Bulgaria și Ungaria, 2019--octombrie 2026; test 2025--2026.}}%
{\T{Quantile MLP (pinball loss) + hour-picking rules; tabular Q-learning; five seeds; NPV.}{MLP cuantilic (pierdere \emph{pinball}) + reguli de alegere a orelor; Q-learning tabelar; cinci seeds; VAN.}}%
{\T{MLP rule \NbestShare\,\% of perfect foresight; agent about as much as fixed hours; \NbestYear\ EUR/year vs \Nneed\ needed.}{Regula MLP: \NbestShare\,\% din previziunea perfectă; agentul, aproximativ cît orele fixe; \NbestYear\ EUR/an față de \Nneed\ necesari.}}%
{\T{Trade on the network's forecast; do not invest on arbitrage alone at 500,000 EUR CAPEX.}{Tranzacționați pe baza prognozei rețelei; nu investiți numai pentru arbitraj la un CAPEX de 500.000 EUR.}}
\begin{center}
\includegraphics[width=0.64\linewidth]{fig_headline_\lang}
\src{\T{our computation on energy-charts.info prices, January 2025 -- October 2026.}{calcule proprii pe prețurile energy-charts.info, ianuarie 2025 -- octombrie 2026.}}
\end{center}
\begin{minipage}[t]{0.32\linewidth}\centering{\color{MainBlue}\Huge\bfseries \NbestShare\,\%}\\ \small \T{of the perfect-foresight profit captured by the quantile MLP}{din profitul previziunii perfecte, captat de MLP-ul cuantilic}\end{minipage}\hfill
\begin{minipage}[t]{0.32\linewidth}\centering{\color{MainBlue}\Huge\bfseries \NbestPay}\\ \small \T{years payback, best realistic tool}{ani, perioada de recuperare pentru cel mai bun instrument realist}\end{minipage}\hfill
\begin{minipage}[t]{0.32\linewidth}\centering{\color{IDAred}\Huge\bfseries \NbestNPV}\\ \small \T{EUR NPV at 500,000 EUR CAPEX}{EUR, VAN la un CAPEX de 500.000 EUR}\end{minipage}
\end{document}
